% Generated draft: requires mathematical and source review.
\documentclass[11pt]{article}
\usepackage[margin=25mm]{geometry}
\usepackage{fontspec}
\setmainfont{Latin Modern Roman}
\usepackage{amsmath,amssymb,mathtools,unicode-math}
\setmathfont{Latin Modern Math}
\usepackage{enumitem,needspace,xurl,hyperref}
\setlength{\parindent}{0pt}
\setlength{\parskip}{5pt}
\setlength{\emergencystretch}{3em}
\newcommand{\N}{\mathbb{N}}
\newcommand{\Z}{\mathbb{Z}}
\newcommand{\Q}{\mathbb{Q}}
\newcommand{\R}{\mathbb{R}}
\newcommand{\C}{\mathbb{C}}
\newcommand{\F}{\mathbb{F}}
\newcommand{\A}{\mathbb{A}}
\newcommand{\PP}{\mathbb{P}}
\newcommand{\E}{\mathbb{E}}
\newcommand{\eps}{\varepsilon}
\newcommand{\dd}{\,\mathrm{d}}
\begin{document}
\section*{UnsolvedMath ID: 30006559}
\textbf{Source problem code:} catalogue entry
\subsection{Definitions and mathematical statement}
Let \(n\in\Z_{\geq1}\) and let \(p\) be real with \(1\leq p<\infty\). Define the metric torus \(T_p^n\) to be \([0,1)^n\), with coordinates interpreted modulo \(1\), equipped with the distance
\[
 d_p(a,b)=\left(\sum_{i=1}^{n}\min\{|a_i-b_i|,\,1-|a_i-b_i|\}^{p}\right)^{1/p}.
\]
A subset \(A\) of a metric space \((X,d)\) is equilateral if there is a number \(r>0\) such that \(d(x,y)=r\) for every two distinct points \(x,y\in A\). Write
\[
 e(X)=\sup\{\#A:A\subseteq X\text{ is a finite equilateral set}\}.
\]
Every equilateral subset of \(T_p^n\) is finite, so this definition also captures the source's supremum over all equilateral subsets.

The first question is whether
\[
 \forall p\in[1,\infty)\ \exists C_p>0\ \forall n\in\Z_{\geq1},\qquad e(T_p^n)\leq C_p n.
\]
Here \(C_p\) may depend on \(p\) but must be independent of \(n\). The further question is whether
\[
 \forall p\in[1,\infty)\ \forall n\in\Z_{\geq1},\qquad
 2n+1\text{ prime}\ \Longrightarrow\ e(T_p^n)=2n+1.
\]
Both questions concern finite \(p\).

The supplied exposition states the lower bounds \(e(T_p^n)\geq2n\), and \(e(T_p^n)\geq2n+1\) when \(2n+1\) is prime. It describes the first construction as \(\{\pm\alpha e_1,\ldots,\pm\alpha e_n\}\), interpreted modulo \(1\), with \(e_i\) the coordinate vectors and \(\alpha\) chosen appropriately; its value is not supplied. The prime lower bound alone does not establish the proposed equality.

The supplied context also reports polynomial upper bounds for fixed finite \(p\). Its flattened expression \(C_p n^{2+2/\lfloor p\rfloor}\) is a reading of the transcription whose exponent typography requires checking against the report. The comparison for \(p=\infty\) is transcribed with the string \(3n\); whether a superscript was lost there remains unresolved. Neither contextual bound is a hypothesis of the questions. These supplied contextual claims have not been independently verified here.
\subsection{Short English statement}
For each fixed finite \(p\geq1\), must every equilateral set in the \(n\)-dimensional torus with its coordinatewise circular \(\ell_p\) metric have at most \(C_p n\) points, uniformly in \(n\)? When \(2n+1\) is prime, is the maximum exactly \(2n+1\)? The supplied record labels these questions open; this is a statement-only draft requiring mathematical review.
\subsection{Sources}
\begin{itemize}
\item Local UnsolvedMath registry, permanent ID \(30006559\), source problem code \texttt{catalogue entry}, title: Equilateral Sets in Finite-\(p\) Metric Tori. This draft uses the supplied complete source record.
\item \emph{Combinatorics} (2026), Oberwolfach Reports, Report 1/2026, pp.~5--79, DOI \texttt{10.4171/owr/2026/1}: \url{https://doi.org/10.4171/owr/2026/1}. This report is supplied as the source of the workshop problem discussion.
\item Alon, \emph{Equilateral sets in the torus}: \url{https://web.math.princeton.edu/~nalon/PDFS/equitor.pdf}. This is a supplied literature source; no further bibliographic details were provided.
\item Attribution: the supplied exposition credits Kravitz and Steinerberger with suggesting consideration of this family of metric spaces. The extracted text places the name Noga Alon near the relevant definition, but also interleaves Hans Yu and an unrelated incidence-geometry question. The structured proposer field is empty, so a definitive proposer attribution cannot be recovered from this record alone.
\item Supplied literature assessment dated exactly \texttt{2026-08-22}: the questions are labelled open; finite-\(p\) upper bounds are described as polynomial and far from linear, and the prime lower bound is said not to establish equality. This is supplied dated curation, not an independently verified literature assessment.
\item The registry labels the statement \texttt{corrected\_verified} and says that the official report was checked and the notation repaired. These are supplied source-verification assertions, not verification performed for this draft.
\end{itemize}
\end{document}
